\documentclass[11pt,a4paper]{scrartcl}
\usepackage[T1]{fontenc}
\usepackage[utf8]{inputenc}
\usepackage[ngerman]{babel}
\usepackage{lmodern}
\usepackage{microtype}
\usepackage[margin=2.5cm]{geometry}
\usepackage{amsmath,amssymb,amsthm}
\usepackage{graphicx,booktabs}
\usepackage{xcolor}
\usepackage{csquotes}
\usepackage[backend=biber,style=numeric,sorting=none]{biblatex}
\addbibresource{literatur.bib}
\usepackage[hidelinks]{hyperref}
\hypersetup{pdftitle={Die Schnorr-Identifikation},pdfauthor={Fabian Grebert}}
%
% Für die Abgabe auf \entwurffalse setzen. Zuvor alle TODOs abarbeiten.
\newif\ifentwurf
\entwurftrue
\newcommand{\todo}[1]{\ifentwurf\par\smallskip\noindent
  {\color{blue!60!black}\textbf{Schreibhinweis:} #1}\par\smallskip\fi}
\newcommand{\Z}{\mathbb{Z}}
\newcommand{\G}{\mathbb{G}}
\newtheorem{definition}{Definition}[section]
\newtheorem{satz}{Satz}[section]
\theoremstyle{remark}
\newtheorem{beispiel}{Beispiel}[section]

\title{Die Schnorr-Identifikation}
\subtitle{Mathematische Grundlagen, Protokoll und praktische Einordnung}
\author{Fabian Grebert}
\date{\today} % Vor Abgabe gegebenenfalls festes Datum eintragen.
%
\begin{document}
\maketitle
\begin{abstract}
\todo{Zum Schluss schreiben: Problem, Vorgehen und wichtigste Ergebnisse
in einem kurzen Absatz zusammenfassen. Keine neuen Argumente einführen.}
\end{abstract}
\tableofcontents

% !TeX root = ../main.tex

\section{Einleitung und Problemstellung}\label{sec:einleitung}

Die zuverlässige Authentifizierung von Benutzern und
technischen Systemen stellt eine zentrale Herausforderung
des Identity and Access Managements (IAM) dar.
Bevor einem Benutzer der Zugriff auf schützenswerte
Ressourcen gewährt werden kann, muss zunächst überprüft
werden, ob dieser tatsächlich über die erforderlichen
Authentifizierungsmerkmale verfügt.

Klassische Authentifizierungsverfahren basieren häufig
auf Passwörtern, die als gemeinsames Geheimnis zwischen
einem Benutzer und einem authentifizierenden System
dienen. Solche Verfahren können jedoch verschiedene
Sicherheitsrisiken mit sich bringen. Dazu zählen
insbesondere schwache oder wiederverwendete Passwörter,
Phishing sowie die Kompromittierung serverseitig
gespeicherter Passwort-Hashes.

Asymmetrische kryptografische Verfahren verfolgen
einen anderen Ansatz. Anstelle eines gemeinsamen
Geheimnisses wird ein Schlüsselpaar verwendet, das
aus einem privaten und einem öffentlichen Schlüssel
besteht. Während der private Schlüssel ausschließlich
dem jeweiligen Schlüsselinhaber zur Verfügung stehen
soll, kann der öffentliche Schlüssel anderen
Teilnehmern zugänglich gemacht werden.

Hieraus ergibt sich eine grundlegende kryptografische
Problemstellung: Wie kann eine Partei nachweisen,
dass sie den zu einem öffentlichen Schlüssel
gehörenden geheimen Schlüssel kennt, ohne diesen
gegenüber einer anderen Partei offenzulegen?


\subsection{Problemstellung}

Für die Betrachtung dieser Problemstellung werden
zwei beteiligte Parteien unterschieden:

\begin{description}

    \item[Prover (beweisende Partei):]
    Der Prover besitzt einen geheimen Schlüssel und
    möchte nachweisen, dass er diesen kennt.

    \item[Verifier (prüfende Partei):]
    Der Verifier kennt den zugehörigen öffentlichen
    Schlüssel und möchte überprüfen, ob der Prover
    tatsächlich über das entsprechende geheime
    Wissen verfügt.

\end{description}

Eine naheliegende, jedoch unsichere Möglichkeit
bestünde darin, den geheimen Schlüssel unmittelbar
an den Verifier zu übertragen. Dieser könnte
anschließend überprüfen, ob der übermittelte
Schlüssel zum bekannten öffentlichen Schlüssel
gehört.

Ein solcher Ansatz würde allerdings die Geheimhaltung
des privaten Schlüssels aufheben. Der Verifier
oder ein Angreifer, der Zugriff auf den übertragenen
Schlüssel erlangt, könnte diesen für zukünftige
Authentifizierungsvorgänge verwenden.

Ein geeignetes kryptografisches Verfahren muss
daher einen Nachweis der Schlüsselkenntnis
ermöglichen, ohne dass der geheime Schlüssel
selbst übertragen werden muss.

Darüber hinaus sollte ein aufgezeichneter
Authentifizierungsvorgang nicht ohne Weiteres
für eine spätere Authentifizierung wiederverwendet
werden können.

Für die praktische Anwendung ist außerdem eine
vertrauenswürdige Zuordnung zwischen dem
öffentlichen Schlüssel und der jeweiligen
Benutzeridentität erforderlich.

Der mathematische Nachweis der Kenntnis eines
privaten Schlüssels bestätigt zunächst lediglich,
dass eine Partei über das entsprechende geheime
Wissen verfügt. Ob dieser Schlüssel tatsächlich
einer bestimmten Person zugeordnet werden kann,
hängt von der zuvor erfolgten Registrierung
und Identitätsprüfung ab.

Die kryptografische Identifikation ist somit
ein Bestandteil eines umfassenderen
Authentifizierungs- und Identitätsmanagementsystems.


\subsection{Das Schnorr-Identifikationsprotokoll}

Eine Lösung für die beschriebene Problemstellung
stellt das Schnorr-Identifikationsprotokoll dar.

Das Verfahren geht auf den deutschen Mathematiker
und Kryptografen Claus-Peter Schnorr zurück.
In seiner Veröffentlichung \emph{Efficient Signature
Generation by Smart Cards} aus dem Jahr 1991
beschreibt Schnorr effiziente Verfahren zur
kryptografischen Identifikation und zur Erzeugung
digitaler Signaturen \cite{schnorr1991}.

Die mathematische Grundlage des Verfahrens bildet
das diskrete Logarithmusproblem in geeigneten
zyklischen Gruppen.

Dabei wird aus einem geheimen Exponenten $x$
und einem öffentlich bekannten Gruppenelement $g$
der öffentliche Schlüssel
\[
    y=g^x
\]
berechnet.

Während die Berechnung des öffentlichen Schlüssels
aus dem geheimen Exponenten effizient möglich ist,
gilt die umgekehrte Berechnung des geheimen
Exponenten aus dem öffentlichen Schlüssel bei
geeigneter Wahl der mathematischen Gruppe und
ihrer Parameter als rechnerisch schwierig.

Das Schnorr-Identifikationsprotokoll nutzt diese
Eigenschaft, um einen interaktiven Nachweis der
Kenntnis des geheimen Exponenten zu ermöglichen.

Der Nachweis erfolgt durch den Austausch von
drei aufeinanderfolgenden Nachrichten:

\begin{enumerate}
    \item Commitment,
    \item Challenge,
    \item Response.
\end{enumerate}

Der Verifier überprüft anhand einer mathematischen
Gleichung, ob die empfangene Response mit dem
öffentlichen Schlüssel und den zuvor ausgetauschten
Werten übereinstimmt.

Der geheime Schlüssel selbst wird während der
Protokollausführung nicht übertragen.

Eine besondere Eigenschaft des Schnorr-Protokolls
besteht darin, dass sich die Kenntnis des
geheimen Schlüssels nachweisen lässt, ohne
einem ehrlichen Verifier zusätzliche Informationen
über diesen Schlüssel offenzulegen.

Diese Eigenschaft wird als Honest-Verifier
Zero Knowledge bezeichnet und im weiteren
Verlauf der Arbeit näher untersucht.


\subsection{Abgrenzung zu anderen kryptografischen Verfahren}

Das Schnorr-Identifikationsprotokoll muss von
anderen Verfahren der asymmetrischen Kryptografie
unterschieden werden.

Insbesondere verfolgen Verschlüsselung, digitale
Signaturen und kryptografische Identifikation
unterschiedliche Schutzziele.

\paragraph{Verschlüsselung.}

Das primäre Ziel einer Verschlüsselung besteht
darin, die Vertraulichkeit von Informationen
sicherzustellen.

Eine Nachricht wird mithilfe eines kryptografischen
Verfahrens so transformiert, dass ihr ursprünglicher
Inhalt nur von berechtigten Parteien wiederhergestellt
werden kann.

Das Schnorr-Identifikationsprotokoll dient hingegen
nicht der Verschlüsselung von Nachrichten.
Es ermöglicht den Nachweis kryptografischen Wissens,
ohne die Vertraulichkeit der übertragenen
Kommunikationsinhalte selbst sicherzustellen.

\paragraph{Digitale Signaturen.}

Digitale Signaturen ermöglichen die Überprüfung
der Authentizität und Integrität einer Nachricht.

Hierzu erzeugt der Inhaber eines privaten Schlüssels
eine Signatur für eine bestimmte Nachricht.
Andere Teilnehmer können diese Signatur mithilfe
des zugehörigen öffentlichen Schlüssels überprüfen.

Im Unterschied dazu dient das ursprüngliche
Schnorr-Identifikationsprotokoll dem interaktiven
Nachweis der Kenntnis eines geheimen Schlüssels
gegenüber einem Verifier.

Eine unmittelbare Bindung des Nachweises an eine
zu signierende Nachricht ist im ursprünglichen
Identifikationsprotokoll nicht vorgesehen.

Allerdings lassen sich aus dem Schnorr-Protokoll
durch geeignete Erweiterungen digitale
Signaturverfahren ableiten.

Insbesondere die Fiat-Shamir-Transformation
ermöglicht es, die interaktive Challenge durch
einen kryptografischen Hashwert zu ersetzen.
Dieser Zusammenhang bildet eine Grundlage
für Schnorr-Signaturen.

\paragraph{Kryptografische Identifikation.}

Das Ziel eines kryptografischen
Identifikationsverfahrens besteht darin,
einen Teilnehmer anhand eines ihm zugeordneten
kryptografischen Geheimnisses zu authentifizieren.

Das Schnorr-Protokoll verfolgt dieses Ziel,
indem es die Kenntnis eines geheimen
diskreten Logarithmus nachweist.

Die Identifikation einer konkreten Person
setzt dabei weiterhin voraus, dass der
verwendete öffentliche Schlüssel zuvor
vertrauenswürdig mit ihrer Identität
verknüpft wurde.

Tabelle~\ref{tab:kryptografische-verfahren}
fasst die unterschiedlichen Zielsetzungen
zusammen.

\begin{table}[htbp]
    \centering
    \small
    \renewcommand{\arraystretch}{1.3}
    \begin{tabular}{p{3.5cm}p{8cm}}
        \toprule
        Verfahren & Primäres Ziel\\
        \midrule

        Verschlüsselung
            & Schutz der Vertraulichkeit
              einer Nachricht\\

        Digitale Signatur
            & Nachweis der Authentizität
              und Integrität einer Nachricht\\

        Schnorr-Identifikation
            & Nachweis der Kenntnis eines
              geheimen Schlüssels ohne dessen
              Offenlegung\\

        \bottomrule
    \end{tabular}
    \caption{Abgrenzung ausgewählter
    kryptografischer Verfahren.
    Eigene Darstellung.}
    \label{tab:kryptografische-verfahren}
\end{table}


\subsection{Zielsetzung und Forschungsfrage}

Ziel dieser Arbeit ist es, die mathematischen
Grundlagen und die Funktionsweise des
Schnorr-Identifikationsprotokolls zu untersuchen
sowie dessen Sicherheit und praktische
Einsetzbarkeit im Kontext des Identity and
Access Managements zu bewerten.

Dabei soll insbesondere nachvollzogen werden,
wie ein Teilnehmer die Kenntnis eines
geheimen Schlüssels nachweisen kann, ohne
diesen gegenüber der prüfenden Partei
offenzulegen.

Aus dieser Zielsetzung ergibt sich folgende
zentrale Forschungsfrage:

\begin{quote}
    \textit{Wie ermöglicht das
    Schnorr-Identifikationsprotokoll den Nachweis
    der Kenntnis eines geheimen Schlüssels ohne
    dessen Offenlegung, und unter welchen
    mathematischen und technischen Voraussetzungen
    eignet sich das Verfahren für eine sichere
    und effiziente Authentifizierung?}
\end{quote}

Zur Beantwortung dieser Forschungsfrage werden
folgende Teilfragen untersucht:

\begin{enumerate}

    \item Auf welchen mathematischen Grundlagen
    basiert das Schnorr-Identifikationsprotokoll?

    \item Wie erfolgt der Nachweis der Kenntnis
    eines geheimen Schlüssels innerhalb des
    Drei-Schritt-Protokolls?

    \item Welche kryptografischen
    Sicherheitseigenschaften besitzt das Verfahren,
    und unter welchen Voraussetzungen gelten diese?

    \item Welcher rechnerische Aufwand ist für
    die reguläre Protokollausführung und für
    mögliche Angriffe erforderlich?

    \item Welche Vorteile, Nachteile und
    praktischen Einsatzbedingungen ergeben
    sich für die Verwendung des Verfahrens
    in Authentifizierungssystemen?

\end{enumerate}

Die Untersuchung erfolgt anhand einer theoretischen
Betrachtung der mathematischen Grundlagen,
der Analyse des Protokollablaufs sowie einer
literaturgestützten Bewertung seiner
Sicherheitseigenschaften.

Zur Veranschaulichung werden numerische Beispiele
mit kleinen Gruppenparametern verwendet.
Diese dienen ausschließlich der Nachvollziehbarkeit
der mathematischen Zusammenhänge und stellen
keine kryptografisch sicheren Parametersätze dar.

Ergänzend werden theoretische Aufwandsabschätzungen
für ausgewählte Angriffsverfahren vorgenommen.

Eine experimentelle Untersuchung konkreter
Implementierungen oder eine vollständige
Sicherheitsanalyse sämtlicher möglicher
Angreifermodelle ist nicht Gegenstand dieser
Arbeit.


\subsection{Aufbau der Arbeit}

Die vorliegende Arbeit gliedert sich in
sechs Abschnitte.

Nach der Einführung in die Problemstellung
und der Formulierung der Forschungsfrage
werden in Abschnitt~\ref{sec:grundlagen}
die mathematischen Grundlagen erläutert.

Hierzu zählen insbesondere Gruppen,
zyklische Gruppen, modulare Arithmetik
und das diskrete Logarithmusproblem.

Abschnitt~\ref{sec:protokoll} beschreibt
anschließend das Schnorr-Identifikationsprotokoll.

Neben der Schlüsselerzeugung und dem Ablauf
des Drei-Schritt-Protokolls werden dessen
Korrektheit sowie die wesentlichen
kryptografischen Sicherheitseigenschaften
untersucht.

In Abschnitt~\ref{sec:praxis} erfolgt
eine Betrachtung des praktischen Bezugs
und des rechnerischen Aufwands.

Dabei werden insbesondere generische
Verfahren zur Berechnung diskreter
Logarithmen, theoretische Angriffszeiten
und ausgewählte Anwendungsfälle betrachtet.

Abschnitt~\ref{sec:bewertung} untersucht
die Vor- und Nachteile des Verfahrens
und leitet daraus geeignete
Einsatzbedingungen ab.

Hierzu wird das Schnorr-Protokoll auch
mit anderen Authentifizierungsmechanismen
verglichen.

Abschließend werden in
Abschnitt~\ref{sec:fazit} die gewonnenen
Erkenntnisse zusammengeführt, die
Forschungsfrage beantwortet und die
Grenzen der vorliegenden Untersuchung
aufgezeigt.

% !TeX root = ../main.tex
\section{Mathematische Grundlagen}\label{sec:grundlagen}
\subsection{Gruppen und zyklische Gruppen}
\todo{Gruppe mit Verknüpfung definieren: Abgeschlossenheit, Assoziativität,
 neutrales Element und inverse Elemente. Ein kleines Beispiel geben.
 Danach endliche zyklische Gruppe, Generator und Gruppenordnung erklären.}
% Beispiel für eine Definitionsumgebung:
% \begin{definition}[Zyklische Gruppe]
% Eine Gruppe ...
% \end{definition}
\subsection{Modulare Arithmetik und Untergruppen}
\todo{Rechnen modulo einer Zahl erklären. Für das spätere Beispiel eine
 Untergruppe von $\Z_p^*$ mit primer Ordnung $q$ und $q\mid(p-1)$ einführen.
 Gruppenelemente werden modulo $p$, Exponenten modulo $q$ berechnet.}
\subsection{Das diskrete Logarithmusproblem}
\todo{Für $\G=\langle g\rangle$ und $y=g^x$ das Problem der Bestimmung
 von $x$ definieren. Leichte Vorwärtsberechnung und angenommene Schwierigkeit
 der Umkehrung unterscheiden. Die Schwierigkeit hängt von Gruppe und Parametern ab.}
\subsection{Verwendete Notation}
\begin{center}
\begin{tabular}{ll}
\toprule
Symbol & Bedeutung \\
\midrule
$\G$, $q$, $g$ & Zyklische Gruppe, prime Ordnung, Generator \\
$x\in\Z_q$, $y=g^x$ & Geheimer und öffentlicher Schlüssel \\
$r\in\Z_q$, $t=g^r$ & Frischer Zufallswert und Commitment \\
$c\in\Z_q$, $s\in\Z_q$ & Challenge und Response \\
\bottomrule
\end{tabular}
\end{center}
\todo{Notation im gesamten Paper konsistent halten; abweichende Vorzeichen
 in Quellen ausdrücklich übersetzen. Alle später benötigten Begriffe hier erklären.}


% !TeX root = ../main.tex

\section{Die Schnorr-Identifikation}\label{sec:protokoll}

Das Schnorr-Identifikationsprotokoll ist ein interaktives
kryptografisches Verfahren, mit dem eine Partei die Kenntnis
eines geheimen Schlüssels nachweisen kann, ohne diesen
gegenüber einer anderen Partei offenzulegen.

Das Verfahren basiert auf dem in
Abschnitt~\ref{sec:grundlagen} beschriebenen diskreten
Logarithmusproblem. Es handelt sich um ein sogenanntes
Challenge-Response-Verfahren, bei dem die Authentifizierung
durch den Austausch von drei Nachrichten erfolgt.

Die beweisende Partei wird im Folgenden als \emph{Prover}
und die prüfende Partei als \emph{Verifier} bezeichnet.


\subsection{Parameter und Schlüsselerzeugung}

Die Grundlage des Schnorr-Identifikationsprotokolls bildet
eine endliche zyklische Gruppe
\[
    \G=\langle g\rangle
\]
mit primer Ordnung $q$ und einem Generator $g$.

Für die Realisierung in einer multiplikativen Untergruppe
von $\Z_p^*$ werden zwei Primzahlen $p$ und $q$ gewählt,
für die
\[
    q\mid(p-1)
\]
gilt. Der Generator $g$ besitzt dabei die Ordnung $q$,
sodass
\[
    g^q\equiv 1\pmod p
    \qquad\text{und}\qquad
    g\neq 1
\]
gelten.

Die Parameter $p$, $q$ und $g$ sind öffentlich bekannt
und müssen nicht geheim gehalten werden. Sie sind jedoch
so zu wählen, dass das diskrete Logarithmusproblem in
der verwendeten Gruppe praktisch nicht effizient lösbar
ist.

\paragraph{Schlüsselerzeugung.}

Der Prover erzeugt seinen geheimen Schlüssel $x$, indem
er einen Wert gleichverteilt aus den von null verschiedenen
Elementen von $\Z_q$ auswählt:
\[
    x\xleftarrow{\$}\Z_q\setminus\{0\}.
\]

Anschließend berechnet er den zugehörigen öffentlichen
Schlüssel
\[
    y=g^x.
\]

Bei der konkreten Realisierung in einer Untergruppe von
$\Z_p^*$ erfolgt die Berechnung durch
\[
    y=g^x\bmod p.
\]

Damit besteht das Schlüsselpaar aus
\begin{align*}
    x &\quad\text{(geheimer Schlüssel)},\\
    y &\quad\text{(öffentlicher Schlüssel)}.
\end{align*}

Der geheime Schlüssel verbleibt ausschließlich beim Prover.
Der öffentliche Schlüssel wird dem Verifier zur Verfügung
gestellt und kann beispielsweise während einer Registrierung
einer bestimmten Benutzeridentität zugeordnet werden.

Diese Zuordnung muss vertrauenswürdig erfolgen. Das
Schnorr-Identifikationsprotokoll weist zunächst lediglich
die Kenntnis des zu einem öffentlichen Schlüssel gehörenden
diskreten Logarithmus nach. Eine Aussage über die tatsächliche
Identität einer Person ist nur möglich, wenn der öffentliche
Schlüssel zuvor zuverlässig mit dieser Identität verknüpft
wurde.

Die öffentlichen Parameter und der öffentliche Schlüssel
müssen vor ihrer Verwendung auf Gültigkeit überprüft werden.
Insbesondere muss sichergestellt werden, dass $y$ ein
zulässiges Element der vorgesehenen Untergruppe ist.


\subsection{Ablauf des Drei-Schritt-Protokolls}

Das Schnorr-Identifikationsprotokoll besteht aus drei
aufeinanderfolgenden Nachrichten:

\begin{enumerate}
    \item Commitment des Provers,
    \item Challenge des Verifiers,
    \item Response des Provers.
\end{enumerate}

Ein Protokoll mit dieser Struktur und den später beschriebenen
Sicherheitseigenschaften wird als Sigma-Protokoll
($\Sigma$-Protokoll) bezeichnet.

Zu Beginn der Protokollausführung kennen beide Parteien
die öffentlichen Parameter und den öffentlichen Schlüssel
$y$. Der geheime Schlüssel $x$ ist ausschließlich dem
Prover bekannt.

\subsubsection*{Schritt 1: Commitment}

Der Prover wählt zunächst einen frischen Zufallswert
$r$ gleichverteilt aus $\Z_q$:
\[
    r\xleftarrow{\$}\Z_q.
\]

Dieser temporäre Zufallswert wird auch als Nonce bezeichnet.
Anschließend berechnet der Prover das Commitment
\[
    t=g^r.
\]

Für die Realisierung in $\Z_p^*$ gilt entsprechend
\[
    t=g^r\bmod p.
\]

Der Prover sendet ausschließlich das Commitment $t$ an
den Verifier. Der zugrunde liegende Zufallswert $r$
wird nicht übertragen und muss geheim bleiben.

Das Commitment legt den Prover auf einen bestimmten
Zufallswert fest, bevor ihm die Challenge des Verifiers
bekannt ist.

\subsubsection*{Schritt 2: Challenge}

Nach Erhalt des Commitments überprüft der Verifier dessen
Gültigkeit und wählt anschließend eine Challenge $c$
gleichverteilt aus $\Z_q$:
\[
    c\xleftarrow{\$}\Z_q.
\]

Die Challenge wird an den Prover übertragen.

Entscheidend ist, dass die Challenge erst nach Erhalt
des Commitments gewählt wird und für den Prover im Voraus
nicht vorhersagbar ist.

Dadurch soll verhindert werden, dass ein Teilnehmer ohne
Kenntnis des geheimen Schlüssels bereits vor Beginn
der Prüfung ein passendes Transkript vorbereitet.

\subsubsection*{Schritt 3: Response}

Nach Erhalt der Challenge berechnet der Prover unter
Verwendung seines geheimen Schlüssels $x$ und des zuvor
gewählten Zufallswertes $r$ die Response
\[
    s=(r+cx)\bmod q.
\]

Anschließend sendet der Prover den Wert $s$ an den Verifier.

Der geheime Schlüssel $x$ sowie der temporäre Zufallswert
$r$ werden während des gesamten Protokollablaufs nicht
übertragen.

Die drei Nachrichten lassen sich somit durch folgende
Gleichungen zusammenfassen:
\begin{align}
    t &= g^r
       && \text{Commitment},\\
    c &\xleftarrow{\$}\Z_q
       && \text{Challenge},\\
    s &= (r+cx)\bmod q
       && \text{Response}.
\end{align}

\subsubsection*{Verifikation}

Nach Erhalt der Response besitzt der Verifier sämtliche
für die Prüfung benötigten öffentlichen Werte.

Er überprüft, ob die erhaltenen Werte gültig sind und
ob folgende Gleichung erfüllt ist:
\begin{equation}\label{eq:pruefung}
    g^s\stackrel{?}{=}t\,y^c.
\end{equation}

Im Fall einer multiplikativen Untergruppe von $\Z_p^*$
lautet die Prüfgleichung
\[
    g^s\stackrel{?}{\equiv}t\,y^c\pmod p.
\]

Sind die empfangenen Werte gültig und ist die Gleichung
erfüllt, akzeptiert der Verifier den Nachweis. Andernfalls
wird die Authentifizierung abgelehnt.

\subsubsection*{Darstellung des Nachrichtenflusses}

Der vollständige Nachrichtenfluss zwischen Prover und
Verifier ist in Abbildung~\ref{fig:schnorr-ablauf}
dargestellt.

\begin{figure}[htbp]
    \centering
    \renewcommand{\arraystretch}{1.5}
    \begin{tabular}{ccc}
        \toprule
        \textbf{Prover} &
        &
        \textbf{Verifier}\\
        \midrule
        Kenntnis von $x$ &
        &
        Kenntnis von $y=g^x$\\
        $r\xleftarrow{\$}\Z_q$ &
        & \\
        $t=g^r$ &
        $\xrightarrow{\quad t\quad}$ &
        \\
        &
        $\xleftarrow{\quad c\quad}$ &
        $c\xleftarrow{\$}\Z_q$\\
        $s=(r+cx)\bmod q$ &
        $\xrightarrow{\quad s\quad}$ &
        \\
        &
        &
        $g^s\stackrel{?}{=}t\,y^c$\\
        \bottomrule
    \end{tabular}
    \caption{Nachrichtenfluss des
    Schnorr-Identifikationsprotokolls}
    \label{fig:schnorr-ablauf}
\end{figure}

Die während einer Protokollausführung übertragenen
Nachrichten bilden gemeinsam das sogenannte Transkript
\[
    (t,c,s).
\]

Ein Transkript wird als akzeptierend bezeichnet, wenn
seine Werte gültig sind und die
Prüfgleichung~\eqref{eq:pruefung} erfüllt ist.


\subsection{Korrektheit}

Die Korrektheit des Schnorr-Identifikationsprotokolls
ergibt sich aus den algebraischen Eigenschaften der
verwendeten zyklischen Gruppe.

Es wird angenommen, dass beide Parteien das Protokoll
ordnungsgemäß ausführen und der Prover den geheimen
Schlüssel $x$ kennt.

Aus der Definition des öffentlichen Schlüssels
\[
    y=g^x
\]
sowie des Commitments
\[
    t=g^r
\]
und der Response
\[
    s=(r+cx)\bmod q
\]
folgt aufgrund der Gruppenordnung $q$:
\begin{align}
    g^s
        &=g^{r+cx}\\
        &=g^r\cdot g^{cx}\\
        &=g^r\cdot(g^x)^c\\
        &=t\cdot y^c.
\end{align}

Somit ist die Prüfgleichung
\[
    g^s=t\cdot y^c
\]
für einen ehrlich erzeugten Protokollablauf erfüllt.

Die Reduktion der Response modulo $q$ verändert das
resultierende Gruppenelement nicht, da $g^q=e$ gilt.

Damit ist gezeigt, dass ein Prover, der den geheimen
Schlüssel kennt und das Protokoll korrekt ausführt,
die Prüfung eines ehrlichen Verifiers erfolgreich
bestehen kann.

Diese Eigenschaft wird als Vollständigkeit
(\emph{Completeness}) bezeichnet.


\subsection{Numerisches Beispiel}

Zur Veranschaulichung des Protokollablaufs werden die
bereits in Abschnitt~\ref{sec:grundlagen} eingeführten
Parameter verwendet:
\[
    p=23,\qquad q=11,\qquad g=2.
\]

Die Gruppe $\G=\langle 2\rangle$ besitzt modulo $23$
die Ordnung $11$.

\paragraph{Schlüsselerzeugung.}

Der Prover wählt den geheimen Schlüssel
\[
    x=7.
\]

Der zugehörige öffentliche Schlüssel ergibt sich aus
\begin{align*}
    y&=g^x\bmod p\\
     &=2^7\bmod 23\\
     &=13.
\end{align*}

\paragraph{Commitment.}

Für die Protokollausführung wird beispielhaft der
Zufallswert
\[
    r=4
\]
gewählt.

Das Commitment lautet damit
\begin{align*}
    t&=g^r\bmod p\\
     &=2^4\bmod 23\\
     &=16.
\end{align*}

Der Prover sendet den Wert $t=16$ an den Verifier.

\paragraph{Challenge.}

Der Verifier wählt beispielhaft die Challenge
\[
    c=3
\]
und sendet diese an den Prover.

\paragraph{Response.}

Der Prover berechnet die Response
\begin{align*}
    s&=(r+cx)\bmod q\\
     &=(4+3\cdot 7)\bmod 11\\
     &=25\bmod 11\\
     &=3.
\end{align*}

Der berechnete Wert $s=3$ wird an den Verifier übertragen.

\paragraph{Verifikation.}

Der Verifier überprüft die Gleichung
\[
    g^s\stackrel{?}{\equiv}t\,y^c\pmod p.
\]

Für die linke Seite ergibt sich
\[
    g^s=2^3\bmod 23=8.
\]

Die rechte Seite liefert
\begin{align*}
    t\,y^c\bmod p
        &=16\cdot 13^3\bmod 23\\
        &=35152\bmod 23\\
        &=8.
\end{align*}

Da beide Seiten übereinstimmen, akzeptiert der Verifier
den Nachweis:
\[
    8=8.
\]

Die während der Protokollausführung verwendeten Werte
sind in Tabelle~\ref{tab:schnorr-beispiel} zusammengefasst.

\begin{table}[htbp]
    \centering
    \begin{tabular}{lrl}
        \toprule
        Symbol & Wert & Bedeutung\\
        \midrule
        $p$ & 23 & Primzahl als Modulus\\
        $q$ & 11 & Gruppenordnung\\
        $g$ & 2  & Generator\\
        $x$ & 7  & Geheimer Schlüssel\\
        $y$ & 13 & Öffentlicher Schlüssel\\
        $r$ & 4  & Temporärer Zufallswert\\
        $t$ & 16 & Commitment\\
        $c$ & 3  & Challenge\\
        $s$ & 3  & Response\\
        \bottomrule
    \end{tabular}
    \caption{Beispiel einer erfolgreichen
    Schnorr-Identifikation}
    \label{tab:schnorr-beispiel}
\end{table}

Das resultierende akzeptierende Transkript lautet
\[
    (t,c,s)=(16,3,3).
\]

Der geheime Schlüssel $x$ wurde während der
Protokollausführung nicht direkt an den Verifier übertragen.

Aufgrund der bewusst kleinen Parameter könnte der geheime
Schlüssel in diesem Beispiel allerdings durch Berechnung
des diskreten Logarithmus leicht bestimmt werden.
Das Beispiel dient daher ausschließlich der Veranschaulichung
und stellt keine kryptografisch sichere Implementierung dar.


\subsection{Sicherheitsidee und Grenzen}

Das Schnorr-Identifikationsprotokoll besitzt drei
wesentliche Eigenschaften, die für Sigma-Protokolle
charakteristisch sind:

\begin{itemize}
    \item Vollständigkeit (\emph{Completeness}),
    \item spezielle Wissensextraktion
          (\emph{Special Soundness}),
    \item Zero Knowledge gegenüber einem ehrlichen
          Verifier (\emph{Honest-Verifier Zero Knowledge}).
\end{itemize}

Diese Eigenschaften beschreiben unterschiedliche Aspekte
der Protokollsicherheit und müssen voneinander
unterschieden werden.

\subsubsection*{Vollständigkeit}

Die Vollständigkeit besagt, dass ein ehrlicher Prover,
der den geheimen Schlüssel kennt, einen ehrlichen
Verifier von dieser Kenntnis überzeugen kann.

Wie im Abschnitt zur Korrektheit gezeigt wurde, gilt
für eine ordnungsgemäße Protokollausführung stets
\[
    g^s=t\,y^c.
\]

Unter der Voraussetzung gültiger Parameter und einer
korrekten Ausführung akzeptiert der Verifier daher
den Nachweis mit Wahrscheinlichkeit $1$.

Die Vollständigkeit ist somit eine Aussage über das
Verhalten ehrlicher Teilnehmer und noch kein Nachweis
dafür, dass sich ein Angreifer nicht erfolgreich
authentifizieren kann.


\subsubsection*{Spezielle Wissensextraktion}

Die spezielle Wissensextraktion beschreibt eine
Eigenschaft, die für den Nachweis der Kenntnis des
geheimen Schlüssels wesentlich ist.

Angenommen, es liegen zwei akzeptierende Transkripte
\[
    (t,c_1,s_1)
    \qquad\text{und}\qquad
    (t,c_2,s_2)
\]
mit demselben Commitment $t$, jedoch unterschiedlichen
Challenges
\[
    c_1\neq c_2
\]
vor.

Da beide Transkripte akzeptierend sind, gilt
\begin{align}
    g^{s_1}&=t\,y^{c_1},\\
    g^{s_2}&=t\,y^{c_2}.
\end{align}

Durch Multiplikation der ersten Gleichung mit dem
Inversen der zweiten Gleichung ergibt sich
\[
    g^{s_1-s_2}=y^{c_1-c_2}.
\]

Mit $y=g^x$ folgt
\[
    g^{s_1-s_2}=g^{x(c_1-c_2)}.
\]

Da $g$ die Ordnung $q$ besitzt, müssen die Exponenten
modulo $q$ übereinstimmen:
\[
    s_1-s_2\equiv x(c_1-c_2)\pmod q.
\]

Weil $q$ prim ist und $c_1\neq c_2$ gilt, besitzt
$c_1-c_2$ ein multiplikatives Inverses modulo $q$.

Durch Multiplikation mit diesem Inversen erhält man
\begin{equation}\label{eq:schnorr-extraktion}
    x=(s_1-s_2)(c_1-c_2)^{-1}\bmod q.
\end{equation}

Aus zwei akzeptierenden Transkripten mit identischem
Commitment und unterschiedlichen Challenges kann somit
der geheime Schlüssel effizient berechnet werden.

Diese Eigenschaft wird in einem theoretischen
Sicherheitsbeweis durch einen sogenannten
Wissensextraktor genutzt.

Der Wissensextraktor betrachtet eine Protokollausführung
und versetzt den untersuchten Prover gedanklich in den
Zustand unmittelbar nach dem Senden des Commitments
zurück. Anschließend wird eine andere Challenge
verwendet.

Kann der Prover mit hinreichend hoher Wahrscheinlichkeit
gültige Antworten auf unterschiedliche Challenges
für dasselbe Commitment erzeugen, lässt sich daraus
der zum öffentlichen Schlüssel gehörende diskrete
Logarithmus extrahieren.

Dieses Verfahren wird als \emph{Rewinding} bezeichnet.
Es handelt sich um eine Technik des Sicherheitsbeweises
und nicht um einen regulären Schritt des eigentlichen
Identifikationsprotokolls.

Die spezielle Wissensextraktion bildet damit eine
Grundlage für die Interpretation des Schnorr-Protokolls
als Wissensnachweis.

Für eine gleichverteilt aus $\Z_q$ gewählte Challenge
beträgt der sogenannte Wissensfehler des Sigma-Protokolls
\[
    \varepsilon=\frac{1}{q}.
\]

Ein Teilnehmer kann beispielsweise im Voraus eine
Challenge erraten und ein passendes Transkript
konstruieren. Unter der Annahme einer unabhängigen
und gleichverteilten Challenge beträgt die
Erfolgswahrscheinlichkeit einer solchen Strategie
genau $1/q$.

Ein Prover, dessen Erfolgswahrscheinlichkeit wesentlich
über diesem Wissensfehler liegt, kann unter den
Voraussetzungen des entsprechenden Extraktionsarguments
zur Berechnung eines gültigen geheimen Schlüssels
verwendet werden.

Bei ausreichend großer Gruppenordnung und unter der
Annahme der Schwierigkeit des diskreten Logarithmusproblems
ergibt sich daraus die Sicherheitsidee des
Identifikationsverfahrens.


\subsubsection*{Honest-Verifier Zero Knowledge}

Die Zero-Knowledge-Eigenschaft beschreibt, dass ein
Nachweis keine zusätzlichen Informationen über das
geheime Wissen des Provers offenlegen soll.

Für das Schnorr-Identifikationsprotokoll lässt sich
diese Eigenschaft zunächst für einen ehrlichen Verifier
zeigen.

Ein ehrlicher Verifier hält sich an die
Protokollspezifikation und wählt seine Challenge
insbesondere unabhängig und gleichverteilt aus
dem vorgesehenen Wertebereich.

Die Beweisidee besteht darin, einen Simulator zu
konstruieren, der ein akzeptierendes Transkript
erzeugen kann, ohne den geheimen Schlüssel $x$
zu kennen.

Der Simulator kennt lediglich die öffentlichen Parameter
und den öffentlichen Schlüssel $y$.

Er führt folgende Schritte aus:
\begin{enumerate}
    \item Wähle eine Challenge $c$ gleichverteilt
    aus $\Z_q$:
    \[
        c\xleftarrow{\$}\Z_q.
    \]

    \item Wähle eine Response $s$ unabhängig und
    gleichverteilt aus $\Z_q$:
    \[
        s\xleftarrow{\$}\Z_q.
    \]

    \item Berechne das Commitment rückwärts aus
    der Prüfgleichung:
    \[
        t=g^s\cdot y^{-c}.
    \]
\end{enumerate}

Der Ausdruck $y^{-c}$ bezeichnet dabei die entsprechende
inverse Potenz des Gruppenelements $y$.

Der Simulator gibt anschließend das Transkript
\[
    (t,c,s)
\]
aus.

Dieses Transkript erfüllt die Prüfgleichung, denn
\begin{align*}
    t\,y^c
        &=(g^s\cdot y^{-c})\cdot y^c\\
        &=g^s.
\end{align*}

Es handelt sich somit um ein akzeptierendes Transkript,
obwohl der Simulator den geheimen Schlüssel nicht
kennt.

Darüber hinaus besitzen die simulierten Transkripte
dieselbe Wahrscheinlichkeitsverteilung wie die
Transkripte einer ehrlichen Protokollausführung.

In einer regulären Protokollausführung gilt für jede
fest gewählte Challenge $c$
\[
    s=(r+cx)\bmod q.
\]

Da $r$ gleichverteilt aus $\Z_q$ gewählt wird, ist
auch $s$ gleichverteilt in $\Z_q$. Die Addition des
festen Wertes $cx$ verändert diese Gleichverteilung
nicht.

Das zugehörige Commitment lässt sich durch
\[
    t=g^r
     =g^{s-cx}
     =g^s\cdot y^{-c}
\]
darstellen.

Damit stimmen die Verteilungen der simulierten und
der ehrlich erzeugten Transkripte exakt überein.

Dies wird als perfekte Honest-Verifier-Zero-Knowledge-
Eigenschaft bezeichnet.

Der Verifier erhält aus einer ehrlichen
Protokollausführung folglich keine zusätzliche
Information über den geheimen Schlüssel, die nicht
auch aus den öffentlichen Informationen simuliert
werden könnte.

Dabei ist zu beachten, dass der öffentliche Schlüssel
selbst bereits den geheimen Schlüssel mathematisch
festlegt. Die praktische Geheimhaltung des Schlüssels
beruht weiterhin auf der Schwierigkeit des diskreten
Logarithmusproblems.

\paragraph{Abgrenzung zur Wissenseigenschaft.}

Die Möglichkeit, ein akzeptierendes Transkript ohne
Kenntnis des geheimen Schlüssels zu simulieren, steht
nicht im Widerspruch zur Wissenseigenschaft des
Protokolls.

Der Simulator darf zunächst die Challenge und die
Response auswählen und erst anschließend das passende
Commitment berechnen.

Während einer tatsächlichen Protokollausführung muss
der Prover das Commitment jedoch bereits übertragen,
bevor die Challenge des Verifiers bekannt ist.

Die zeitliche Reihenfolge der Nachrichten ist daher
ein wesentlicher Bestandteil der Protokollsicherheit.

Ein einzelnes akzeptierendes Transkript stellt deshalb
für sich allein keinen unabhängig überprüfbaren Beweis
dar, dass eine bestimmte Person tatsächlich den geheimen
Schlüssel kennt.


\subsubsection*{Bedeutung frischer Zufallswerte}

Die Sicherheit des Schnorr-Identifikationsprotokolls
setzt voraus, dass für jede Protokollausführung ein
neuer, unabhängiger und kryptografisch sicherer
Zufallswert $r$ erzeugt wird.

Wird derselbe Zufallswert in zwei Ausführungen verwendet,
so entsteht dasselbe Commitment
\[
    t=g^r.
\]

Erhält der Prover dabei zwei unterschiedliche Challenges
$c_1$ und $c_2$, ergeben sich die Responses
\begin{align}
    s_1&=(r+c_1x)\bmod q,\\
    s_2&=(r+c_2x)\bmod q.
\end{align}

Durch Subtraktion der beiden Gleichungen fällt der
Zufallswert $r$ weg:
\[
    s_1-s_2\equiv(c_1-c_2)x\pmod q.
\]

Der geheime Schlüssel kann daher unmittelbar mithilfe
von Gleichung~\eqref{eq:schnorr-extraktion} berechnet
werden.

Die Wiederverwendung eines Zufallswertes kann somit
zur vollständigen Kompromittierung des langfristigen
geheimen Schlüssels führen.

\paragraph{Numerisches Beispiel einer
Nonce-Wiederverwendung.}

Es werden erneut die Beispielparameter
\[
    p=23,\qquad q=11,\qquad g=2,\qquad x=7
\]
verwendet.

Angenommen, der Prover verwendet in zwei verschiedenen
Protokollausführungen denselben Zufallswert
\[
    r=4.
\]

Das Commitment lautet in beiden Ausführungen
\[
    t=2^4\bmod 23=16.
\]

Für die erste Challenge $c_1=3$ erhält man
\begin{align*}
    s_1&=(4+3\cdot 7)\bmod 11\\
       &=3.
\end{align*}

Für die zweite Challenge $c_2=5$ ergibt sich
\begin{align*}
    s_2&=(4+5\cdot 7)\bmod 11\\
       &=6.
\end{align*}

Damit liegen die beiden akzeptierenden Transkripte
\[
    (16,3,3)
    \qquad\text{und}\qquad
    (16,5,6)
\]
vor.

Der geheime Schlüssel kann nun durch
\begin{align*}
    x
        &=(s_1-s_2)(c_1-c_2)^{-1}\bmod q\\
        &=(3-6)(3-5)^{-1}\bmod 11\\
        &=(-3)(-2)^{-1}\bmod 11
\end{align*}
berechnet werden.

Da
\[
    (-2)^{-1}\equiv 5\pmod{11}
\]
gilt, ergibt sich
\begin{align*}
    x&=(-3)\cdot 5\bmod 11\\
     &=-15\bmod 11\\
     &=7.
\end{align*}

Der ursprüngliche geheime Schlüssel wurde damit
vollständig rekonstruiert.

Dieses Beispiel verdeutlicht, dass die Sicherheit
des Protokolls nicht ausschließlich von der Schwierigkeit
des diskreten Logarithmusproblems abhängt, sondern
auch von der Qualität und korrekten Verwendung der
erzeugten Zufallswerte.

Auch die Offenlegung oder Vorhersagbarkeit eines
einzelnen Zufallswertes kann den geheimen Schlüssel
gefährden. Für eine von null verschiedene Challenge
$c$ gilt bei bekanntem $r$
\[
    x=(s-r)c^{-1}\bmod q.
\]

Der temporäre Zufallswert muss daher sowohl während
als auch nach seiner Verwendung vor unbefugtem
Zugriff geschützt werden. Insbesondere darf ein
bereits verwendetes Commitment nicht für die
Beantwortung weiterer Challenges wiederverwendet
werden.


\subsubsection*{Angreifermodell und weitere
Sicherheitsgrenzen}

Die beschriebenen mathematischen Eigenschaften gelten
unter bestimmten Voraussetzungen. Sie erlauben keine
pauschale Aussage über die Sicherheit jeder praktischen
Implementierung.

Insbesondere sind folgende Aspekte zu berücksichtigen:

\begin{description}

    \item[Unvorhersehbare Challenges:]
    Die Challenge muss unabhängig vom Prover und
    mithilfe eines kryptografisch sicheren
    Zufallszahlengenerators erzeugt werden.
    Ist die Challenge vorhersehbar, kann ein Angreifer
    ein passendes Commitment vorbereiten, ohne den
    geheimen Schlüssel zu kennen.

    \item[Validierung der Gruppenparameter:]
    Die verwendeten Parameter müssen die notwendigen
    mathematischen Voraussetzungen erfüllen.
    Dazu zählen insbesondere die Primzahleigenschaften
    von $p$ und $q$, die Teilbarkeitsbedingung
    $q\mid(p-1)$ sowie die korrekte Ordnung des
    Generators $g$.

    \item[Validierung der Gruppenelemente:]
    Öffentliche Schlüssel und Commitments müssen
    auf ihre Zugehörigkeit zur vorgesehenen Untergruppe
    überprüft werden.

    Für ein Element $u$ in der multiplikativen
    Untergruppe müssen insbesondere
    \[
        1\leq u\leq p-1
    \]
    und
    \[
        u^q\equiv 1\pmod p
    \]
    gelten.

    Für den öffentlichen Schlüssel ist zusätzlich
    $y\neq 1$ erforderlich, da der zugehörige
    diskrete Logarithmus andernfalls öffentlich
    bekannt wäre.

    \item[Replay-Angriffe:]
    Ein aufgezeichnetes Transkript kann nicht ohne
    Weiteres für eine spätere Authentifizierung
    verwendet werden, da der Verifier eine neue,
    unabhängige Challenge erzeugt.

    Stimmt die neue Challenge mit der aufgezeichneten
    Challenge überein, ist ein Replay jedoch grundsätzlich
    möglich. Bei gleichverteilter Auswahl aus $\Z_q$
    beträgt die Wahrscheinlichkeit hierfür $1/q$.
    Die Challenge muss daher einen ausreichend
    großen Wertebereich besitzen.

    \item[Relay-Angriffe:]
    Eine frische Challenge verhindert nicht zwangsläufig
    die Weiterleitung einer laufenden Authentifizierung.
    Ein Angreifer könnte Nachrichten zwischen einem
    legitimen Prover und einem Verifier weiterleiten,
    ohne den geheimen Schlüssel selbst zu kennen.

    Zusätzliche Mechanismen zur Bindung der
    Authentifizierung an den Kommunikationskanal,
    die Gegenstelle oder den Anwendungskontext
    können daher erforderlich sein.

    \item[Manipulierte Verifier:]
    Die hier gezeigte Zero-Knowledge-Eigenschaft
    bezieht sich auf einen ehrlichen Verifier.
    Daraus folgt nicht unmittelbar eine entsprechende
    Sicherheit gegenüber beliebig vom Protokoll
    abweichenden Verifiern.

    Weitergehende Aussagen erfordern ein präzises
    Sicherheitsmodell und entsprechende Beweise.

    \item[Implementierungssicherheit:]
    Fehler bei der Implementierung, unsichere
    Zufallszahlengeneratoren, Seitenkanalangriffe
    oder ein unzureichender Schutz des geheimen
    Schlüssels können die Sicherheit des Verfahrens
    beeinträchtigen.

    \item[Quantencomputer:]
    Die langfristige Sicherheit des Verfahrens
    hängt von der Schwierigkeit des diskreten
    Logarithmusproblems ab.
    Ein hinreichend leistungsfähiger, fehlertoleranter
    Quantencomputer könnte dieses Problem mithilfe
    des Shor-Algorithmus effizient lösen.

    Das klassische Schnorr-Identifikationsprotokoll
    ist daher kein postquantenkryptografisches
    Verfahren.

\end{description}


\subsection{Abgrenzung zu nichtinteraktiven Verfahren}

Das in diesem Abschnitt beschriebene
Schnorr-Identifikationsprotokoll setzt eine aktive
Interaktion zwischen Prover und Verifier voraus.

Eine bekannte Erweiterung besteht darin, die zufällige
Challenge des Verifiers durch einen aus öffentlichen
Informationen berechneten kryptografischen Hashwert
zu ersetzen. Dieser Ansatz wird als
Fiat-Shamir-Transformation bezeichnet.

Abstrakt kann eine solche Challenge beispielsweise
durch
\[
    c=H(g,y,t,\text{Kontext})\bmod q
\]
dargestellt werden.

Der Prover kann dadurch die Challenge selbst berechnen,
ohne zuvor eine Nachricht des Verifiers erhalten
zu müssen.

Die konkrete Definition und Kodierung der Hash-Eingaben
sowie die Abbildung des Hashwertes in den vorgesehenen
Challenge-Bereich sind dabei sicherheitsrelevant.

Diese Transformation ermöglicht unter geeigneten
Annahmen die Konstruktion nichtinteraktiver
Wissensnachweise und bildet außerdem eine Grundlage
für Schnorr-Signaturverfahren.

Bei einer digitalen Signatur werden zusätzlich die
zu signierende Nachricht und gegebenenfalls weitere
Kontextinformationen in die Berechnung der Challenge
einbezogen.

Die Sicherheit einer solchen Erweiterung muss jedoch
gesondert untersucht werden. Die Eigenschaften des
interaktiven Protokolls können nicht ohne weitere
Annahmen unmittelbar auf jede nichtinteraktive
Variante übertragen werden.


\subsection{Zusammenfassung}

Das Schnorr-Identifikationsprotokoll ermöglicht den
Nachweis der Kenntnis eines geheimen diskreten
Logarithmus, ohne diesen direkt an den Verifier
zu übertragen.

Die Kombination aus Commitment, Challenge und
Response gewährleistet bei einer korrekten
Protokollausführung die Vollständigkeit des Verfahrens.

Die spezielle Wissensextraktion zeigt, dass aus
zwei akzeptierenden Transkripten mit identischem
Commitment und unterschiedlichen Challenges der
geheime Schlüssel berechnet werden kann.

Gleichzeitig lässt sich durch Simulation nachweisen,
dass eine ehrliche Protokollausführung keine
zusätzlichen Informationen über den geheimen
Schlüssel offenlegt.

Die Sicherheit des Verfahrens beruht somit auf dem
Zusammenspiel der mathematischen Eigenschaften des
Protokolls, der Schwierigkeit des diskreten
Logarithmusproblems und der korrekten praktischen
Implementierung.


% !TeX root = ../main.tex

\section{Praxisbezug und Aufwandsabschätzung}\label{sec:praxis}

Nachdem in den vorherigen Abschnitten die mathematischen
Grundlagen und der Ablauf des Schnorr-Identifikationsprotokolls
beschrieben wurden, wird im Folgenden dessen praktische
Bedeutung untersucht.

Hierzu werden zunächst der Aufwand möglicher Angriffe
und die Auswirkungen unterschiedlicher Gruppenordnungen
betrachtet. Anschließend werden konkrete Anwendungsfälle
sowie verwandte kryptografische Verfahren vorgestellt.


\subsection{Ein kleines Rechenbeispiel}

Das numerische Beispiel aus Abschnitt~\ref{sec:protokoll}
veranschaulicht die Funktionsweise des Schnorr-Protokolls
anhand der Parameter
\[
    p=23,\qquad q=11,\qquad g=2.
\]

Für den geheimen Schlüssel $x=7$, den Zufallswert $r=4$
und die Challenge $c=3$ ergibt sich das akzeptierende
Transkript
\[
    (t,c,s)=(16,3,3).
\]

Der zugehörige öffentliche Schlüssel lautet $y=13$.

Obwohl das Protokoll mathematisch korrekt ausgeführt
wird, bieten diese Parameter keine praktische
kryptografische Sicherheit.

Ein Angreifer, der lediglich die öffentlichen Werte
$p$, $q$, $g$ und $y$ kennt, kann sämtliche möglichen
privaten Schlüssel ausprobieren.

Die folgende Tabelle zeigt die Berechnung der Potenzen
des Generators $g=2$ modulo $23$.

\begin{table}[htbp]
    \centering
    \begin{tabular}{rr}
        \toprule
        Exponent $k$ & $g^k\bmod p$\\
        \midrule
        0 & 1\\
        1 & 2\\
        2 & 4\\
        3 & 8\\
        4 & 16\\
        5 & 9\\
        6 & 18\\
        7 & 13\\
        8 & 3\\
        9 & 6\\
        10 & 12\\
        \bottomrule
    \end{tabular}
    \caption{Vollständige Suche nach dem diskreten
    Logarithmus in einer kleinen zyklischen Gruppe}
    \label{tab:bruteforce-beispiel}
\end{table}

Da der öffentliche Schlüssel den Wert $y=13$ besitzt,
kann aus Tabelle~\ref{tab:bruteforce-beispiel}
unmittelbar der geheime Schlüssel
\[
    x=7
\]
abgelesen werden.

Ein Angreifer könnte sich anschließend mit diesem
Schlüssel gegenüber einem Verifier authentifizieren.

Das Beispiel zeigt, dass die mathematische Korrektheit
eines kryptografischen Protokolls allein keine
ausreichende Sicherheit garantiert.

Die Sicherheit hängt wesentlich von der Größe und
Struktur der verwendeten Gruppe sowie von der Qualität
der Implementierung ab.


\subsection{Aufwand eines Angriffs}

Für die Bewertung der Sicherheit des
Schnorr-Identifikationsprotokolls müssen unterschiedliche
Angriffsziele voneinander unterschieden werden.

Ein Angreifer kann beispielsweise versuchen,

\begin{enumerate}
    \item den langfristigen geheimen Schlüssel $x$
    aus dem öffentlichen Schlüssel $y$ zu berechnen,

    \item einen Verifier erfolgreich zu täuschen,
    ohne den geheimen Schlüssel zu kennen,

    \item Schwachstellen der Implementierung
    auszunutzen, etwa eine fehlerhafte Erzeugung
    temporärer Zufallswerte.
\end{enumerate}

Die ersten beiden Angriffsziele werden im Folgenden
näher untersucht.

Angriffe auf die Implementierung, insbesondere die
Wiederverwendung von Zufallswerten, wurden bereits in
Abschnitt~\ref{sec:protokoll} betrachtet.


\subsubsection{Schlüsselrückgewinnung durch Brute Force}

Ein grundlegender Angriff auf das Schnorr-Protokoll
besteht darin, den diskreten Logarithmus des öffentlichen
Schlüssels zu berechnen.

Gegeben sind die öffentlichen Parameter
$\G$, $q$, $g$ sowie der öffentliche Schlüssel
\[
    y=g^x.
\]

Gesucht wird der geheime Exponent
\[
    x\in\Z_q\setminus\{0\}.
\]

Bei einer vollständigen Suche, auch als Brute-Force-Angriff
bezeichnet, werden sämtliche möglichen Exponenten
systematisch überprüft.

Hierzu wird für jeden Kandidaten $k$ geprüft, ob
\[
    g^k=y
\]
gilt.

Sobald die Gleichung erfüllt ist, wurde der gesuchte
diskrete Logarithmus gefunden.

Die einzelnen Potenzen müssen dabei nicht vollständig
neu berechnet werden. Stattdessen können sie schrittweise
durch Multiplikation mit dem Generator erzeugt werden:
\[
    g^{k+1}=g^k\cdot g.
\]

Bei einem gleichverteilt gewählten geheimen Schlüssel
sind durchschnittlich etwa
\[
    \frac{q}{2}
\]
Versuche erforderlich.

Im schlechtesten Fall müssen alle $q-1$ möglichen
von null verschiedenen Exponenten überprüft werden.

Die asymptotische Laufzeit beträgt somit
\[
    \mathcal{O}(q)
\]
Gruppenoperationen.

Da bei der Suche lediglich eine konstante Anzahl
von Gruppenelementen gleichzeitig gespeichert werden
muss, beträgt der zusätzliche Speicherbedarf
\[
    \mathcal{O}(1)
\]
Gruppenelemente.

Eine vollständige Suche ist bei ausreichend großer
Gruppenordnung praktisch nicht durchführbar.
Allerdings existieren effizientere Algorithmen zur
Berechnung diskreter Logarithmen
\cite[Abschnitt 3.6.1]{menezes1996}.


\subsubsection{Baby-Step--Giant-Step}

Der Baby-Step--Giant-Step-Algorithmus ist ein
generisches Verfahren zur Berechnung diskreter
Logarithmen in zyklischen Gruppen.

Im Unterschied zur vollständigen Suche reduziert
der Algorithmus die Anzahl erforderlicher
Gruppenoperationen durch die Verwendung einer
vorberechneten Tabelle.

Hierzu wird zunächst
\[
    m=\left\lceil\sqrt{q}\right\rceil
\]
gewählt.

Jeder gesuchte Exponent $x$ mit $0\leq x<q$
kann in der Form
\[
    x=im+j
\]
mit geeigneten ganzen Zahlen $i,j$ und
$0\leq i,j<m$ dargestellt werden.

Aus
\[
    y=g^{im+j}
\]
folgt durch Umformung
\[
    y\cdot g^{-im}=g^j.
\]

Der Algorithmus berechnet zunächst die sogenannten
Baby Steps
\[
    g^j
    \qquad\text{für}\qquad
    0\leq j<m
\]
und speichert sie in einer Tabelle.

Anschließend werden die Giant Steps
\[
    y\cdot g^{-im}
\]
berechnet und mit den gespeicherten Gruppenelementen
verglichen.

Wird eine Übereinstimmung gefunden, kann der
gesuchte Exponent durch
\[
    x=im+j\bmod q
\]
bestimmt werden.

Der Algorithmus benötigt
\[
    \mathcal{O}(\sqrt{q})
\]
Gruppenoperationen sowie Speicherplatz für
\[
    \mathcal{O}(\sqrt{q})
\]
Gruppenelemente.

Bei Verwendung einer geeigneten Hashtabelle können
die gespeicherten Werte effizient nachgeschlagen
werden. Bei einer sortierten Tabelle entsteht
zusätzlicher Aufwand für die Sortierung und Suche.

Damit ist der Baby-Step--Giant-Step-Algorithmus
deutlich effizienter als eine vollständige Suche.

Sein erheblicher Speicherbedarf begrenzt jedoch
die praktische Anwendbarkeit bei großen Gruppen
\cite[Abschnitt 3.6.2]{menezes1996}.


\subsubsection{Pollard-Rho-Algorithmus}

Ein weiteres generisches Verfahren zur Berechnung
diskreter Logarithmen ist der Pollard-Rho-Algorithmus.

Dieser verwendet eine Folge pseudozufällig erzeugter
Gruppenelemente und sucht nach einer Kollision,
also nach zwei übereinstimmenden Gruppenelementen.

Die mathematische Grundlage des Verfahrens ist
mit dem Geburtstagsparadoxon verwandt.

Für eine Menge mit $q$ Elementen ist unter einer
geeigneten Zufallsannahme bereits nach ungefähr
$\sqrt{q}$ Beobachtungen mit einer Kollision
zu rechnen.

Der Pollard-Rho-Algorithmus benötigt daher im
Erwartungswert
\[
    \mathcal{O}(\sqrt{q})
\]
Gruppenoperationen.

Im Gegensatz zum Baby-Step--Giant-Step-Algorithmus
müssen jedoch nicht sämtliche bisher erzeugten
Gruppenelemente gespeichert werden.

Bei Verwendung eines geeigneten Verfahrens zur
Zyklenerkennung benötigt die klassische serielle
Variante lediglich eine konstante Anzahl von
Gruppenelementen und Skalaren:
\[
    \mathcal{O}(1).
\]

Die tatsächlich benötigte Anzahl an Operationen
ist zufallsabhängig. Ein erfolgloser Durchlauf
kann mit anderen Startwerten wiederholt werden.

Durch die Kombination aus vergleichsweise geringer
Laufzeitkomplexität und geringem Speicherbedarf
ist Pollard-Rho für die Bewertung der Sicherheit
großer zyklischer Gruppen besonders relevant
\cite[Abschnitt 3.6.3]{menezes1996}.


\subsubsection{Vergleich der Angriffsverfahren}

Tabelle~\ref{tab:dlog-algorithmen} fasst die
asymptotischen Eigenschaften der beschriebenen
generischen Angriffsverfahren zusammen.

\begin{table}[htbp]
    \centering
    \begin{tabular}{lll}
        \toprule
        Verfahren & Gruppenoperationen & Speicherbedarf\\
        \midrule
        Brute Force
            & $\mathcal{O}(q)$
            & $\mathcal{O}(1)$\\
        Baby-Step--Giant-Step
            & $\mathcal{O}(\sqrt{q})$
            & $\mathcal{O}(\sqrt{q})$\\
        Pollard-Rho
            & $\mathcal{O}(\sqrt{q})$
            & $\mathcal{O}(1)$\\
        \bottomrule
    \end{tabular}
    \caption{Vergleich generischer Algorithmen zur
    Berechnung diskreter Logarithmen.
    Die Angaben zu Pollard-Rho beziehen sich auf den
    erwarteten Aufwand der klassischen seriellen Variante.
    Der Speicherbedarf wird in Gruppenelementen und
    Skalaren angegeben. Eigene Darstellung nach
    \cite[Abschnitt 3.6]{menezes1996}.}
    \label{tab:dlog-algorithmen}
\end{table}

Aus dem Vergleich ergibt sich, dass die Sicherheit
einer zyklischen Gruppe nicht allein anhand der
Anzahl möglicher privater Schlüssel bewertet
werden darf.

Besitzt die Gruppe eine Ordnung mit ungefähr
$n$ Bit, also
\[
    q\approx 2^n,
\]
so benötigen generische Verfahren wie Pollard-Rho
ungefähr
\[
    \sqrt{q}\approx 2^{n/2}
\]
Gruppenoperationen.

Eine Gruppenordnung von ungefähr $256$ Bit
entspricht daher hinsichtlich solcher generischen
Angriffe einem Aufwand in der Größenordnung
von $2^{128}$ Gruppenoperationen.

Diese Aussage beschreibt ausschließlich die
generische Angriffskomplexität und setzt voraus,
dass keine effizienteren gruppenspezifischen
Angriffe anwendbar sind.


\subsection{Konkrete Zeitabschätzung}

Um die Größenordnung des erforderlichen
Rechenaufwands zu verdeutlichen, wird im Folgenden
eine hypothetische Zeitabschätzung vorgenommen.

Dabei wird angenommen, dass einem Angreifer eine
Rechenleistung von
\[
    R=10^{12}
\]
Gruppenoperationen pro Sekunde zur Verfügung steht.

Diese Rate ist eine bewusst gewählte Modellannahme
und kein gemessener Leistungswert realer Hardware.

Insbesondere hängt die tatsächlich erreichbare Rate
von der verwendeten Gruppe, der Größe ihrer Elemente,
der Implementierung und der eingesetzten Hardware ab.

Für einen generischen Angriff mit einem angenommenen
Aufwand von
\[
    N=\sqrt{q}
\]
Gruppenoperationen ergibt sich die geschätzte
Angriffsdauer zu
\begin{equation}\label{eq:angriffsdauer}
    T=\frac{\sqrt{q}}{R}.
\end{equation}

Konstante Faktoren, zusätzlicher Rechenaufwand und
mögliche Parallelisierungseffekte werden in dieser
vereinfachten Betrachtung nicht berücksichtigt.

\subsubsection*{Beispiel für eine Gruppenordnung von
128 Bit}

Für eine Gruppenordnung von ungefähr
\[
    q\approx 2^{128}
\]
ergibt sich ein generischer Rechenaufwand von
\[
    N\approx 2^{64}
     \approx 1{,}84\cdot 10^{19}
\]
Gruppenoperationen.

Unter der angenommenen Operationsrate beträgt die
Angriffsdauer
\begin{align*}
    T
        &\approx\frac{2^{64}}{10^{12}}\\
        &\approx 1{,}84\cdot 10^7
        \text{ Sekunden}.
\end{align*}

Dies entspricht ungefähr
\[
    T\approx 213{,}5\text{ Tagen}.
\]

Das Beispiel zeigt, dass eine Gruppenordnung von
128 Bit nicht mit einem Sicherheitsniveau von
128 Bit gleichgesetzt werden darf.


\subsubsection*{Beispiel für eine Gruppenordnung von
256 Bit}

Für eine Gruppenordnung von ungefähr
\[
    q\approx 2^{256}
\]
ergibt sich ein generischer Rechenaufwand von
\[
    N\approx 2^{128}
     \approx 3{,}40\cdot 10^{38}
\]
Gruppenoperationen.

Unter derselben angenommenen Operationsrate beträgt
die Angriffsdauer
\begin{align*}
    T
        &\approx\frac{2^{128}}{10^{12}}\\
        &\approx 3{,}40\cdot 10^{26}
        \text{ Sekunden}.
\end{align*}

Unter Verwendung von 365,25 Tagen pro Jahr ergibt
sich daraus eine Dauer von ungefähr
\[
    T\approx 1{,}08\cdot 10^{19}
    \text{ Jahren}.
\]

Diese Größenordnung liegt weit außerhalb praktisch
relevanter Zeiträume.

Tabelle~\ref{tab:angriffszeiten} vergleicht die
geschätzte Dauer für unterschiedliche Bitlängen
der Gruppenordnung.

\begin{table}[htbp]
    \centering
    \begin{tabular}{rrl}
        \toprule
        Bitlänge von $q$
            & Generischer Aufwand
            & Geschätzte Dauer\\
        \midrule
        64 Bit  & $2^{32}$  & 4,29 Millisekunden\\
        128 Bit & $2^{64}$  & 213,5 Tage\\
        160 Bit & $2^{80}$  & $3{,}83\cdot 10^4$ Jahre\\
        192 Bit & $2^{96}$  & $2{,}51\cdot 10^9$ Jahre\\
        256 Bit & $2^{128}$ & $1{,}08\cdot 10^{19}$ Jahre\\
        \bottomrule
    \end{tabular}
    \caption{Theoretische Zeitabschätzung generischer
    Angriffe bei angenommenen $10^{12}$ Gruppenoperationen
    pro Sekunde. Eigene Berechnung gemäß
    Gleichung~\eqref{eq:angriffsdauer}.}
    \label{tab:angriffszeiten}
\end{table}

Die dargestellten Werte sind keine Prognosen für
konkrete Angriffe. Sie dienen ausschließlich der
Veranschaulichung des Zusammenhangs zwischen
Gruppenordnung und generischer Angriffskomplexität.


\subsubsection{Speicherbedarf generischer Angriffe}

Neben dem zeitlichen Aufwand ist auch der
Speicherbedarf eines Angriffs von Bedeutung.

Für eine Gruppe mit einer Ordnung von ungefähr
$2^{256}$ müsste der Baby-Step--Giant-Step-Algorithmus
eine Tabelle mit ungefähr
\[
    2^{128}
\]
Gruppenelementen anlegen.

Wird beispielhaft eine Darstellung von 32 Byte
pro Gruppenelement angenommen, ergibt sich bereits
für diese Elemente ein Speicherbedarf von
\begin{align*}
    S
        &\approx 32\cdot 2^{128}\\
        &=2^{133}\text{ Byte}\\
        &\approx 1{,}09\cdot 10^{40}
        \text{ Byte}.
\end{align*}

Die Speicherung zugehöriger Exponenten sowie
zusätzlicher Tabellenstrukturen würde den Bedarf
weiter erhöhen.

Die angenommene Elementgröße dient lediglich der
Illustration. Bei einer multiplikativen Untergruppe
eines großen endlichen Körpers können Gruppenelemente
wesentlich mehr Speicherplatz benötigen.

Der Pollard-Rho-Algorithmus besitzt hingegen einen
wesentlich geringeren Speicherbedarf und ist
deshalb für generische Sicherheitsabschätzungen
besonders relevant.


\subsubsection{Spezialisierte Angriffe auf endliche Körper}

Die bisher betrachteten Algorithmen sind generisch,
da sie keine besonderen algebraischen Eigenschaften
der verwendeten Gruppe ausnutzen.

Für bestimmte Gruppen existieren jedoch
spezialisierte Verfahren, die effizienter als
generische Algorithmen sein können.

Bei multiplikativen Gruppen endlicher Körper zählen
hierzu insbesondere Index-Calculus-Verfahren und
darauf aufbauende Varianten des Number Field Sieve.

Diese Verfahren nutzen die konkrete mathematische
Darstellung der Gruppenelemente aus.

Für große Primkörper ist der asymptotische Aufwand
des allgemeinen Number Field Sieve zur Berechnung
diskreter Logarithmen subexponentiell in der
Bitlänge des Körpermodulus.

Daher hängt die Sicherheit einer Untergruppe von
$\Z_p^*$ nicht nur von ihrer Gruppenordnung $q$,
sondern auch wesentlich von der Größe und Struktur
des Primzahlmodulus $p$ ab
\cite[Kapitel 15]{galbraith2012}.

Eine große Gruppenordnung allein garantiert
in diesem Fall kein entsprechendes Sicherheitsniveau.

Bei geeigneten elliptischen Kurven sind hingegen
keine klassischen subexponentiellen Algorithmen
zur Lösung des allgemeinen elliptischen diskreten
Logarithmusproblems bekannt.

Für solche Kurven stellen generische Angriffe
wie Pollard-Rho daher eine wesentliche Grundlage
der Sicherheitsabschätzung dar.

Diese Aussage gilt nicht uneingeschränkt für
jede elliptische Kurve. Bestimmte spezielle Kurven
können zusätzliche mathematische Schwächen
aufweisen.

Für die praktische Auswahl kryptografischer
Parameter müssen deshalb sowohl die Gruppenordnung
als auch die Struktur der zugrunde liegenden
Gruppe berücksichtigt werden.


\subsection{Erfolgreiche Täuschung ohne Schlüsselkenntnis}

Neben der vollständigen Rückgewinnung des geheimen
Schlüssels besteht ein weiteres Angriffsziel darin,
eine einzelne Protokollausführung erfolgreich
abzuschließen, ohne den privaten Schlüssel zu kennen.

Wie in Abschnitt~\ref{sec:protokoll} gezeigt wurde,
kann ein Teilnehmer ein akzeptierendes Transkript
für eine im Voraus festgelegte Challenge erzeugen.

Hierzu wählt er eine Challenge $c'$ und eine
Response $s'$ und berechnet
\[
    t'=g^{s'}y^{-c'}.
\]

Sendet er dieses Commitment an einen ehrlichen
Verifier, so kann er die vorbereitete Response
verwenden, falls die tatsächliche Challenge
mit der zuvor gewählten Challenge übereinstimmt.

Für eine unabhängige und gleichverteilt aus
$\Z_q$ gewählte Challenge gilt
\begin{equation}\label{eq:taeuschungswahrscheinlichkeit}
    \Pr[c=c']=\frac{1}{q}.
\end{equation}

Dies ist die Erfolgswahrscheinlichkeit dieser
konkreten Strategie und entspricht dem in
Abschnitt~\ref{sec:protokoll} beschriebenen
Wissensfehler.

Sie darf nicht als allgemeine obere Schranke
für sämtliche denkbaren Angriffe interpretiert
werden.

\paragraph{Beispiel mit kleinen Parametern.}

Für die verwendete Lernumgebung mit $q=11$
beträgt die Erfolgswahrscheinlichkeit
\[
    \Pr[c=c']
        =\frac{1}{11}
        \approx 9{,}09\,\%.
\]

Ein Angreifer hätte mit dieser Strategie
somit bereits bei einer einzelnen Protokollrunde
eine vergleichsweise hohe Erfolgswahrscheinlichkeit.

\paragraph{Beispiel mit realistischen Größenordnungen.}

Bei einer Challenge aus einem Wertebereich mit
$2^{128}$ gleichwahrscheinlichen Möglichkeiten
beträgt die Erfolgswahrscheinlichkeit des
einfachen Ratens
\[
    \Pr[c=c']=2^{-128}
    \approx 2{,}94\cdot 10^{-39}.
\]

Werden mehrere unabhängige Authentifizierungsversuche
zugelassen, steigt die Wahrscheinlichkeit,
mindestens einen erfolgreichen Versuch zu erzielen.

Für $k$ unabhängige Versuche mit einer jeweiligen
Erfolgswahrscheinlichkeit von $1/q$ ergibt sich
\[
    P_k=1-\left(1-\frac{1}{q}\right)^k.
\]

Für $k\ll q$ kann diese Wahrscheinlichkeit durch
\[
    P_k\approx\frac{k}{q}
\]
angenähert werden.

In praktischen Anwendungen können zusätzliche
Maßnahmen, beispielsweise die Begrenzung der
Anzahl von Authentifizierungsversuchen, das Risiko
solcher Online-Angriffe weiter reduzieren.


\subsection{Rechenaufwand des legitimen Protokolls}

Neben dem Aufwand möglicher Angriffe ist auch die
Effizienz des Protokolls für legitime Teilnehmer
von Bedeutung.

Das Schnorr-Identifikationsprotokoll ist
vergleichsweise effizient, da es nur wenige
aufwendige Gruppenoperationen benötigt.

Der Prover führt pro Protokollausführung im
Wesentlichen folgende Berechnungen durch:
\[
    t=g^r
\]
und
\[
    s=(r+cx)\bmod q.
\]

Die Berechnung des Commitments erfordert eine
Gruppenexponentiation.

Die Berechnung der Response besteht hingegen
lediglich aus einer modularen Multiplikation
und einer modularen Addition.

Der Verifier überprüft anschließend
\[
    g^s=t\,y^c.
\]

Eine direkte Implementierung benötigt hierfür
zwei Gruppenexponentiationen und eine
Gruppenmultiplikation.

Durch Verfahren zur simultanen Exponentiation
kann der Rechenaufwand der Verifikation
reduziert werden.

Tabelle~\ref{tab:legitimer-aufwand} fasst die
wesentlichen Berechnungen zusammen.

\begin{table}[htbp]
    \centering
    \begin{tabular}{lll}
        \toprule
        Partei & Berechnung & Aufwand\\
        \midrule
        Prover
            & $t=g^r$
            & Eine Exponentiation\\
        Prover
            & $s=(r+cx)\bmod q$
            & Skalare Arithmetik\\
        Verifier
            & $g^s\stackrel{?}{=}t\,y^c$
            & Zwei Exponentiationen\\
        \bottomrule
    \end{tabular}
    \caption{Wesentliche Berechnungen einer regulären
    Schnorr-Identifikation ohne zusätzliche
    Parameter- und Schlüsselvalidierung.
    Eigene Darstellung.}
    \label{tab:legitimer-aufwand}
\end{table}

Die Anzahl benötigter Gruppenoperationen für
eine einzelne Exponentiation wächst bei Verwendung
geeigneter Algorithmen, beispielsweise Square-and-Multiply,
nur logarithmisch mit der Größe des Exponenten.

Bei einem Exponenten aus $\Z_q$ ergibt sich
ein Aufwand von
\[
    \mathcal{O}(\log q)
\]
Gruppenoperationen.

Damit besteht ein erheblicher Unterschied
zwischen dem Aufwand der legitimen Teilnehmer
und dem Aufwand eines generischen Angreifers:
\[
    \mathcal{O}(\log q)
    \qquad\text{gegenüber}\qquad
    \mathcal{O}(\sqrt{q}).
\]

Diese Asymmetrie ist eine wesentliche Grundlage
der praktischen Einsetzbarkeit des Verfahrens.

Weitere Ausführungen zur Effizienz der verwandten
nichtinteraktiven Variante finden sich in
\cite[Abschnitte 2.4 und 3.4]{hao2017}.


\subsection{Anwendungsfälle und verwandte Verfahren}

Das Schnorr-Identifikationsprotokoll stellt ein
grundlegendes kryptografisches Verfahren dar,
das sowohl zur direkten Identifikation als
auch als Baustein komplexerer Verfahren
eingesetzt werden kann.

Dabei muss zwischen dem ursprünglichen interaktiven
Protokoll und den daraus abgeleiteten Verfahren
unterschieden werden.


\subsubsection{Interaktive Schlüsselauthentifizierung}

Ein grundlegender Anwendungsfall des
Schnorr-Identifikationsprotokolls ist die
Authentifizierung eines Teilnehmers anhand
eines asymmetrischen Schlüsselpaares.

Bei einer Registrierung erzeugt der Teilnehmer
einen privaten Schlüssel $x$ und hinterlegt
den zugehörigen öffentlichen Schlüssel
\[
    y=g^x.
\]

Bei einer späteren Anmeldung weist der Teilnehmer
mithilfe des interaktiven Schnorr-Protokolls
nach, dass er den zum registrierten öffentlichen
Schlüssel gehörenden privaten Schlüssel kennt.

Der private Schlüssel wird dabei nicht an
den authentifizierenden Server übertragen.

Dieser Ansatz unterscheidet sich von einer
klassischen passwortbasierten Authentifizierung,
bei der ein gemeinsames Geheimnis beziehungsweise
eine daraus abgeleitete Information zur
Überprüfung verwendet wird.

Bei einer schlüsselbasierten Authentifizierung
muss der Server hingegen lediglich den
öffentlichen Schlüssel zur Überprüfung
des mathematischen Nachweises speichern.

Die Offenlegung des öffentlichen Schlüssels
ermöglicht unter der Annahme der Schwierigkeit
des diskreten Logarithmusproblems keine
effiziente Berechnung des privaten Schlüssels.

Allerdings sind für eine sichere praktische
Authentifizierung zusätzliche Mechanismen
erforderlich.

Dazu zählen insbesondere eine vertrauenswürdige
Registrierung des öffentlichen Schlüssels, der
Schutz des privaten Schlüssels und die Bindung
der Authentifizierung an die vorgesehene
Gegenstelle.

Die hier beschriebene Anwendung stellt einen
konzeptionellen Anwendungsfall des interaktiven
Protokolls dar und ist nicht mit einem bestimmten
standardisierten Anmeldedienst gleichzusetzen.


\subsubsection{Nichtinteraktive Wissensnachweise}

Eine wesentliche Erweiterung des Schnorr-Protokolls
ist der nichtinteraktive Zero-Knowledge-Nachweis.

Im Gegensatz zur ursprünglichen interaktiven
Variante muss der Verifier keine zufällige
Challenge an den Prover übertragen.

Stattdessen wird die Challenge mittels einer
kryptografischen Hashfunktion aus dem Commitment,
dem öffentlichen Schlüssel und weiteren
Kontextinformationen abgeleitet.

Diese Transformation geht auf das
Fiat-Shamir-Verfahren zurück und wurde bereits
in Abschnitt~\ref{sec:protokoll} vorgestellt.

RFC~8235 beschreibt eine konkrete nichtinteraktive
Variante des Schnorr-Protokolls sowohl für
multiplikative Untergruppen endlicher Körper
als auch für elliptische Kurven.

Das Dokument ist als informatorischer RFC
veröffentlicht und beschreibt unter anderem
Anwendungsfälle sowie sicherheitsrelevante
Implementierungsvoraussetzungen.

Nichtinteraktive Wissensnachweise können als
Baustein kryptografischer Protokolle verwendet
werden, um die Kenntnis eines diskreten
Logarithmus nachzuweisen, ohne den zugrunde
liegenden Exponenten offenzulegen
\cite[Abschnitte 2, 3 und 5]{hao2017}.


\subsubsection{Passwortauthentifizierter Schlüsselaustausch}

Ein konkreter Anwendungsfall des nichtinteraktiven
Schnorr-Wissensnachweises ist das J-PAKE-Protokoll
(\emph{Password-Authenticated Key Exchange by Juggling}).

J-PAKE ermöglicht es zwei Parteien, auf Grundlage
eines gemeinsam bekannten Passworts einen
kryptografischen Sitzungsschlüssel auszuhandeln.

Hierzu tauschen die Teilnehmer unter anderem
Gruppenelemente aus, deren zugehörige Exponenten
sie mithilfe von Schnorr-Wissensnachweisen
nachweisen.

Diese Nachweise dienen dazu, die korrekte
Erzeugung bestimmter Protokollwerte zu überprüfen,
ohne die jeweiligen geheimen Exponenten zu
übertragen.

Das Verfahren wird in RFC~8236 beschrieben
\cite[Abschnitte 1 und 2]{hao2017jpake}.

Der Anwendungsfall verdeutlicht, dass Schnorr-Nachweise
auch innerhalb komplexerer kryptografischer
Protokolle eingesetzt werden können.

Dabei handelt es sich nicht um eine unmittelbare
Ausführung des ursprünglichen interaktiven
Identifikationsprotokolls, sondern um die
Integration seiner nichtinteraktiven Variante
in ein umfassenderes Schlüsselaustauschverfahren.


\subsubsection{Nachweis des Besitzes eines privaten
Schlüssels}

Ein weiterer möglicher Anwendungsfall ist der
sogenannte \emph{Proof of Possession} (PoP)
bei der Registrierung eines öffentlichen
Schlüssels.

Beispielsweise kann eine Zertifizierungsstelle
vor der Ausstellung eines digitalen Zertifikats
einen Nachweis verlangen, dass der Antragsteller
tatsächlich über den zum eingereichten öffentlichen
Schlüssel gehörenden privaten Schlüssel verfügt.

Ein Schnorr-Wissensnachweis kann für einen
solchen Nachweis verwendet werden, sofern er
eindeutig an den jeweiligen Registrierungsvorgang
gebunden ist.

Hierzu können beispielsweise Informationen
über die Zertifizierungsstelle, den Antragsteller
oder den konkreten Vorgang in die Berechnung
der nichtinteraktiven Challenge einbezogen
werden.

RFC~8235 nennt diesen Anwendungsfall ausdrücklich
als mögliche Verwendung des Schnorr-Nachweises
\cite[Abschnitt 5]{hao2017}.

Dadurch wird nicht automatisch die tatsächliche
Identität des Antragstellers nachgewiesen.

Der Nachweis bezieht sich zunächst ausschließlich
auf die Kenntnis des zum öffentlichen Schlüssel
gehörenden privaten Schlüssels.

Die Überprüfung der Identität stellt einen
zusätzlichen organisatorischen und technischen
Prozess dar.


\subsubsection{Schnorr-Signaturen}

Eine weitere wichtige Entwicklung auf Grundlage
des Schnorr-Identifikationsprotokolls sind
digitale Schnorr-Signaturen.

Bei digitalen Signaturen wird der Nachweis
mit einer konkreten Nachricht verknüpft.

In einer vereinfachten Darstellung wird die
Challenge aus dem Commitment, dem öffentlichen
Schlüssel und der zu signierenden Nachricht
$m$ abgeleitet:
\[
    c=H(t,y,m)\bmod q.
\]

Anschließend wird die Response wie beim
interaktiven Protokoll berechnet:
\[
    s=(r+cx)\bmod q.
\]

Eine mögliche abstrakte Darstellung der
Signatur ist
\[
    \sigma=(t,s).
\]

Der Verifier berechnet die Challenge aus
den öffentlichen Informationen erneut und
überprüft die zugehörige Verifikationsgleichung.

Diese Darstellung verdeutlicht lediglich
den mathematischen Zusammenhang zwischen
Schnorr-Identifikation und Schnorr-Signaturen.

Konkrete Signaturstandards definieren zusätzliche
Vorgaben für die Kodierung der Werte,
die Erzeugung temporärer Skalare und die
Verwendung der Hashfunktion.

Ein praktisches Beispiel für den Einsatz
von Schnorr-Signaturen ist Bitcoin.

Mit dem Bitcoin Improvement Proposal 340
(BIP~340) wurde ein Schnorr-Signaturverfahren
auf Grundlage der elliptischen Kurve
secp256k1 spezifiziert
\cite{wuille2020bip340}.

Schnorr-Signaturen werden dort im Rahmen von
Taproot zur Autorisierung entsprechender
Transaktionen eingesetzt.

Dieser Anwendungsfall zeigt die praktische
Bedeutung der mathematischen Prinzipien des
Schnorr-Protokolls, obwohl nicht das
ursprüngliche interaktive Identifikationsverfahren
ausgeführt wird.


\subsubsection{Abgrenzung zu modernen Passkeys}

Auch moderne passwortlose Authentifizierungsverfahren,
wie beispielsweise Passkeys, verwenden asymmetrische
Schlüsselpaare und kryptografische Nachweise
über den Besitz eines privaten Schlüssels.

Dadurch besteht eine konzeptionelle Ähnlichkeit
zum Schnorr-Identifikationsprotokoll.

Passkeys sind jedoch nicht mit dem
Schnorr-Protokoll gleichzusetzen.

Sie beruhen auf den durch FIDO2 und WebAuthn
definierten Authentifizierungsmechanismen und
können unterschiedliche digitale
Signaturverfahren verwenden.

Darüber hinaus werden die Authentifizierungsdaten
an den jeweiligen Webdienst und den
Authentifizierungsvorgang gebunden.

Die Gemeinsamkeit besteht somit vor allem
in der Verwendung asymmetrischer Kryptografie
zur Authentifizierung, ohne den privaten
Schlüssel an den Server übertragen zu müssen.

Eine direkte Verwendung des interaktiven
Schnorr-Protokolls durch Passkeys kann daraus
nicht abgeleitet werden.


\subsection{Zusammenfassung}

Die Untersuchung des Angriffsaufwands zeigt,
dass die Sicherheit des Schnorr-Protokolls
wesentlich von den gewählten kryptografischen
Parametern abhängt.

Während eine vollständige Suche nach dem
geheimen Schlüssel einen Aufwand von
$\mathcal{O}(q)$ Gruppenoperationen besitzt,
reduzieren generische Algorithmen wie
Baby-Step--Giant-Step und Pollard-Rho diesen
Aufwand auf $\mathcal{O}(\sqrt{q})$.

Für bestimmte mathematische Gruppen können
darüber hinaus spezialisierte Angriffe
existieren.

Eine sichere Parameterauswahl muss daher
sowohl die Gruppenordnung als auch die
Struktur der zugrunde liegenden Gruppe
berücksichtigen.

Gleichzeitig benötigen die legitimen Teilnehmer
des Protokolls vergleichsweise wenige
Gruppenoperationen.

Diese Kombination aus effizienter regulärer
Protokollausführung und hohem angenommenem
Angriffsaufwand macht die zugrunde liegenden
mathematischen Prinzipien für praktische
kryptografische Anwendungen interessant.

Die betrachteten Anwendungsfälle verdeutlichen
zudem, dass das Schnorr-Protokoll nicht nur
als eigenständiges Identifikationsverfahren,
sondern insbesondere als Grundlage für
nichtinteraktive Wissensnachweise, komplexere
Schlüsselaustauschprotokolle und digitale
Signaturen von Bedeutung ist.

% !TeX root = ../main.tex
\section{Vor- und Nachteile sowie Einsatzbedingungen}\label{sec:bewertung}
\subsection{Vorteile}
\todo{Unter anderem Rechen- und Kommunikationsaufwand sowie Nachweis ohne
 Offenlegung des Geheimnisses bewerten. Jede Aussage begründen oder belegen.}
\subsection{Nachteile und Voraussetzungen}
\todo{Interaktion, Zufallszahlenerzeugung, Schlüsselzuordnung, sichere
 Implementierung und Grenzen der zugrunde liegenden Annahmen besprechen.}
\subsection{Wann lohnt sich das Verfahren?}
\todo{Aus den vorherigen Punkten konkrete Einsatzbedingungen ableiten.
 Bezug zur Problemstellung herstellen und Vergleichskriterien offenlegen.}

% !TeX root = ../main.tex
\section{Fazit}\label{sec:fazit}
\todo{Leitfrage beantworten, wichtigste Erkenntnisse zusammenführen und
 Grenzen des Papers benennen. Optional kurzer Ausblick. Keine neuen Beweise
 oder unvorbereiteten technischen Begriffe einführen.}

% Zeigt im Template die Startliteratur. Für die Abgabe entfernen und nur
% tatsächlich verwendete Quellen mittels \cite{...} im Text zitieren.
\nocite{schnorr1991,hao2017,wikipedia}
\printbibliography[heading=bibintoc,title={Literaturverzeichnis}]
\end{document}
